\documentclass[11pt]{article}
\usepackage[margin=1in]{geometry}
\usepackage{amsmath, amssymb}
\usepackage{xcolor}
\usepackage{tcolorbox}

\newcommand{\dd}[2]{\frac{d #1}{d #2}}
\newcommand{\ddd}[2]{\frac{d^2 #1}{d #2^2}}

\title{Block on a Rough Surface Under a Triangular Force Pulse\\[6pt]
\large Governing ODE with Static/Kinetic Friction Transitions}
\author{Claude Opus 4.6}
\date{April 2026}

\begin{document}
\maketitle

%% ============================================================
\section{Problem Setup}

\begin{tcolorbox}[colback=blue!5, colframe=blue!50!black, title=Physical System]
A block of mass $m$ sits at rest on a rough horizontal surface.  Starting at $t = 1\;\text{s}$, a horizontal force $F(t)$ is applied as a \textbf{single symmetric triangular pulse} with peak at $t = 3\;\text{s}$.  The peak force $F_{\max}$ exceeds the maximum static friction $f_s = \mu_s m g$.  The pulse ends at $t = 5\;\text{s}$.  We seek $x(t)$, the position of the block.
\end{tcolorbox}

\medskip
\noindent
\textbf{Given parameters:}
\begin{itemize}
    \item Mass: $m$
    \item Coefficients of friction: $\mu_s$ (static), $\mu_k$ (kinetic), with $\mu_k < \mu_s$
    \item Gravitational acceleration: $g$
    \item Peak applied force: $F_{\max} > \mu_s m g$
    \item Pulse timing: starts at $t = 1$, peaks at $t = 3$, ends at $t = 5$ (all in seconds)
\end{itemize}

%% ============================================================
\section{The Applied Force}

The triangular pulse rises linearly from $0$ to $F_{\max}$ over $[1, 3]$, then falls linearly from $F_{\max}$ to $0$ over $[3, 5]$:

\begin{equation}
\boxed{
F(t) =
\begin{cases}
0 & t < 1 \\[4pt]
\displaystyle F_{\max}\,\frac{t - 1}{2} & 1 \le t \le 3 \\[8pt]
\displaystyle F_{\max}\,\frac{5 - t}{2} & 3 < t \le 5 \\[4pt]
0 & t > 5
\end{cases}
}
\label{eq:Ft}
\end{equation}

\noindent
Note that $F(1) = 0$, $F(3) = F_{\max}$, $F(5) = 0$, and the pulse is symmetric about $t = 3$.

%% ============================================================
\section{Friction Model}

The friction force $f(t)$ depends on whether the block is stationary or sliding.

\medskip
\noindent
\textbf{Static friction} (block at rest, $v = 0$):
\begin{equation}
f = F(t) \quad \text{provided } |F(t)| \le \mu_s m g
\label{eq:static}
\end{equation}
Static friction exactly balances the applied force --- the block does not move.  Motion initiates when $F(t)$ first exceeds $\mu_s m g$.

\medskip
\noindent
\textbf{Kinetic friction} (block sliding, $v > 0$):
\begin{equation}
f = \mu_k m g \quad \text{(opposing the direction of motion)}
\label{eq:kinetic}
\end{equation}

\noindent
Since $F(t) \ge 0$ and the block starts from rest, the velocity will always be $v \ge 0$ (the block only moves in the direction of the applied force).  So kinetic friction acts in the $-x$ direction whenever the block is sliding.

%% ============================================================
\section{Critical Times}

\subsection{Onset of motion: $t = t_1$}

The block begins to slide when $F(t)$ first exceeds the static friction limit.  During the rising phase ($1 \le t \le 3$):
\begin{equation}
F(t_1) = \mu_s m g
\qquad\Longrightarrow\qquad
F_{\max}\,\frac{t_1 - 1}{2} = \mu_s m g
\end{equation}

\begin{equation}
\boxed{t_1 = 1 + \frac{2\,\mu_s m g}{F_{\max}}}
\label{eq:t1}
\end{equation}

\noindent
Since $F_{\max} > \mu_s m g$, we have $t_1 < 3$, confirming the block starts moving during the rising phase.  In fact, $1 < t_1 < 3$.

\subsection{Force drops below kinetic friction: $t = t_2$}

Once the block is sliding, friction is only $\mu_k m g < \mu_s m g$.  The net force on the block becomes zero when $F(t) = \mu_k m g$.  This happens on the \emph{falling} side of the pulse ($3 < t \le 5$):
\begin{equation}
F_{\max}\,\frac{5 - t_2}{2} = \mu_k m g
\end{equation}

\begin{equation}
\boxed{t_2 = 5 - \frac{2\,\mu_k m g}{F_{\max}}}
\label{eq:t2}
\end{equation}

\noindent
For $t > t_2$, the applied force is less than kinetic friction, so friction decelerates the block.

\medskip
\noindent
\textit{Note:} By the symmetry of the pulse, there is also a time on the rising side where $F(t) = \mu_k m g$.  But this occurs \emph{before} the block starts moving (since $\mu_k < \mu_s$), so it plays no dynamical role.

%% ============================================================
\section{Phases of Motion}

The dynamics splits into four distinct phases:

\begin{tcolorbox}[colback=green!5, colframe=green!50!black, title=Phase Summary]
\begin{enumerate}
    \item \textbf{Phase I:} $\;0 \le t < t_1$ --- Block at rest.  Static friction balances $F(t)$.
    \item \textbf{Phase II:} $\;t_1 \le t \le 3$ --- Block sliding, force rising.  Net force accelerates block.
    \item \textbf{Phase III:} $\;3 < t \le t_2$ --- Block sliding, force falling but still $> \mu_k mg$.  Block still accelerates (but less).
    \item \textbf{Phase IV:} $\;t > t_2$ --- Applied force $< \mu_k mg$ (or zero).  Friction decelerates block until it stops.
\end{enumerate}
\end{tcolorbox}

%% ============================================================
\section{Governing ODE by Phase}

Newton's second law in the horizontal direction:
\begin{equation}
m\,\ddd{x}{t} = F(t) - f(t)
\label{eq:newton}
\end{equation}

\subsection{Phase I: $0 \le t < t_1$ (at rest)}

\begin{equation}
\boxed{\ddd{x}{t} = 0, \qquad x(t) = x_0, \qquad v(t) = 0}
\end{equation}

\noindent
Nothing happens.  Static friction $f = F(t)$ exactly cancels the applied force.

\subsection{Phase II: $t_1 \le t \le 3$ (sliding, force rising)}

The applied force is $F(t) = F_{\max}(t-1)/2$ and kinetic friction acts:

\begin{equation}
m\,\ddd{x}{t} = F_{\max}\,\frac{t - 1}{2} - \mu_k m g
\end{equation}

\begin{equation}
\boxed{\ddd{x}{t} = \frac{F_{\max}}{2m}(t - 1) - \mu_k g}
\label{eq:phase2}
\end{equation}

\noindent
with initial conditions $x(t_1) = x_0$, $\;\dot{x}(t_1) = 0$.

\medskip
\noindent
\textbf{First integration} (velocity):
\begin{align}
v(t) &= \int_{t_1}^{t} \left[\frac{F_{\max}}{2m}(\tau - 1) - \mu_k g\right] d\tau \\[6pt]
     &= \frac{F_{\max}}{2m}\left[\frac{(\tau-1)^2}{2}\right]_{t_1}^{t} - \mu_k g\,(t - t_1) \\[6pt]
     &= \frac{F_{\max}}{4m}\left[(t-1)^2 - (t_1-1)^2\right] - \mu_k g\,(t - t_1)
\label{eq:v_phase2}
\end{align}

\medskip
\noindent
\textbf{Second integration} (position):
\begin{align}
x(t) &= x_0 + \frac{F_{\max}}{4m}\int_{t_1}^{t}\left[(\tau-1)^2 - (t_1-1)^2\right]d\tau - \mu_k g\int_{t_1}^{t}(\tau - t_1)\,d\tau \\[6pt]
     &= x_0 + \frac{F_{\max}}{4m}\left[\frac{(t-1)^3 - (t_1-1)^3}{3} - (t_1-1)^2(t - t_1)\right] - \frac{\mu_k g}{2}(t - t_1)^2
\label{eq:x_phase2}
\end{align}

\subsection{Phase III: $3 < t \le t_2$ (sliding, force falling)}

The applied force is now $F(t) = F_{\max}(5 - t)/2$:

\begin{equation}
\boxed{\ddd{x}{t} = \frac{F_{\max}}{2m}(5 - t) - \mu_k g}
\label{eq:phase3}
\end{equation}

\noindent
with initial conditions $x(3) = x_3$, $\;\dot{x}(3) = v_3$ obtained by evaluating \eqref{eq:v_phase2} and \eqref{eq:x_phase2} at $t = 3$.

\medskip
\noindent
\textbf{First integration:}
\begin{align}
v(t) &= v_3 + \int_{3}^{t}\left[\frac{F_{\max}}{2m}(5 - \tau) - \mu_k g\right]d\tau \\[6pt]
     &= v_3 + \frac{F_{\max}}{2m}\left[-(5-\tau)^2/2\,\Big|_{\,3}^{\,t}\right] - \mu_k g\,(t - 3) \\[6pt]
     &= v_3 - \frac{F_{\max}}{4m}\left[(5-t)^2 - 4\right] - \mu_k g\,(t - 3)
\label{eq:v_phase3}
\end{align}

\medskip
\noindent
\textbf{Second integration:}
\begin{align}
x(t) &= x_3 + v_3(t - 3) - \frac{F_{\max}}{4m}\left[-\frac{(5-t)^3 - 8}{3} - 4(t-3)\right] - \frac{\mu_k g}{2}(t-3)^2
\end{align}

\noindent
which simplifies to:
\begin{align}
x(t) &= x_3 + v_3(t-3) + \frac{F_{\max}}{4m}\left[\frac{(5-t)^3 - 8}{3} + 4(t-3)\right] - \frac{\mu_k g}{2}(t-3)^2
\label{eq:x_phase3}
\end{align}

\subsection{Phase IV: $t > t_2$ (deceleration to rest)}

For $t_2 < t \le 5$, the applied force $F(t) < \mu_k mg$, and for $t > 5$, $F(t) = 0$.  In both sub-cases the net force opposes motion.

\medskip
\noindent
\textbf{Sub-phase IVa:} $t_2 < t \le 5$ \quad(applied force still present but insufficient)
\begin{equation}
\ddd{x}{t} = \frac{F_{\max}}{2m}(5-t) - \mu_k g \quad < 0
\label{eq:phase4a}
\end{equation}

\noindent
(This is the same ODE as Phase III, but now the RHS is negative --- the block decelerates.)

\medskip
\noindent
\textbf{Sub-phase IVb:} $t > 5$ \quad(no applied force, pure friction braking)
\begin{equation}
\boxed{\ddd{x}{t} = -\mu_k g}
\label{eq:phase4b}
\end{equation}

\noindent
with initial conditions from matching at $t = 5$.  The block decelerates uniformly until:
\begin{equation}
v(t_{\text{stop}}) = 0
\qquad\Longrightarrow\qquad
t_{\text{stop}} = 5 + \frac{v_5}{\mu_k g}
\label{eq:tstop}
\end{equation}
where $v_5 = v(5)$ from Sub-phase IVa (if the block is still moving at $t = 5$).

\medskip
\noindent
\textit{Important:} The block may stop \emph{during} Sub-phase IVa (before the pulse ends) if $v$ reaches zero while $t < 5$.  In that case, one must check whether static friction can hold the block at rest --- i.e., whether $F(t_{\text{stop}}) \le \mu_s mg$ at the instant of stopping.  If so, the block re-sticks and remains at rest.

%% ============================================================
\section{Compact ODE Summary}

Defining the net acceleration $a(t)$ piecewise, the complete equation of motion is:

\begin{tcolorbox}[colback=red!5, colframe=red!50!black, title=Complete Equation of Motion]
\begin{equation}
\boxed{
m\,\ddot{x} =
\begin{cases}
0 & 0 \le t < t_1 \;\text{(static equilibrium)} \\[6pt]
\displaystyle \frac{F_{\max}}{2}(t-1) - \mu_k m g & t_1 \le t \le 3 \;\text{(sliding, force rising)} \\[8pt]
\displaystyle \frac{F_{\max}}{2}(5-t) - \mu_k m g & 3 < t \le 5 \;\text{(sliding, force falling)} \\[8pt]
-\mu_k m g & t > 5 \;\text{(friction braking)}
\end{cases}
}
\label{eq:master}
\end{equation}

\medskip
\noindent
\textbf{Subject to:}
\begin{itemize}
    \item $x(0) = x_0$, \;$\dot{x}(0) = 0$
    \item The constraint $\dot{x}(t) \ge 0$ \;(the block cannot move backward under forward forcing)
    \item The block re-sticks (returns to $\ddot{x} = 0$) whenever $\dot{x} \to 0^+$ and $F(t) \le \mu_s mg$
\end{itemize}
\end{tcolorbox}

%% ============================================================
\section{Dimensionless Form}

\noindent
It is natural to non-dimensionalize using the following scales:

\begin{align}
\hat{t} &= \frac{t - 1}{2}, \qquad \text{so } \hat{t} = 0 \text{ at pulse start, } \hat{t} = 1 \text{ at peak, } \hat{t} = 2 \text{ at pulse end} \\[6pt]
\hat{x} &= \frac{x - x_0}{F_{\max}/(m \cdot 1\,\text{s}^{-2})}, \qquad \hat{F} = \frac{F}{F_{\max}}
\end{align}

\noindent
Define the two key dimensionless ratios:
\begin{equation}
\boxed{
\alpha = \frac{\mu_s m g}{F_{\max}} \quad (\text{static friction ratio}),
\qquad
\beta = \frac{\mu_k m g}{F_{\max}} \quad (\text{kinetic friction ratio})
}
\label{eq:dimless_params}
\end{equation}

\noindent
The problem requires $\alpha < 1$ (the peak exceeds static friction) and $\beta < \alpha$ (kinetic $<$ static).

\medskip
\noindent
In these variables, the triangular pulse is simply:
\begin{equation}
\hat{F}(\hat{t}) =
\begin{cases}
\hat{t} & 0 \le \hat{t} \le 1 \\
2 - \hat{t} & 1 < \hat{t} \le 2
\end{cases}
\end{equation}
and the ODE for $\hat{t} > \hat{t}_1 = \alpha$ becomes:
\begin{equation}
\boxed{\ddd{\hat{x}}{\hat{t}} = \hat{F}(\hat{t}) - \beta}
\label{eq:dimless_ode}
\end{equation}

\noindent
This is the cleanest form of the problem: a piecewise-linear forcing minus a constant friction, with onset at $\hat{t} = \alpha$.

%% ============================================================
\section{Key Observations}

\begin{enumerate}
    \item \textbf{The ODE is piecewise-linear, hence exactly solvable.}  Each phase yields polynomials in $t$ (quadratic for velocity, cubic for position) joined by continuity conditions at the transition times.

    \item \textbf{The gap $\mu_s - \mu_k$ creates an asymmetry.}  The block requires a force $\mu_s mg$ to start moving, but once sliding, friction drops to $\mu_k mg$.  This means the block \emph{cannot} re-stick until the applied force drops below $\mu_s mg$ on the falling side --- and even then, only if the velocity happens to reach zero at that instant.

    \item \textbf{The block always overshoots.}  Because $F_{\max} > \mu_s mg$ and the pulse is symmetric, the block accumulates significant velocity by $t = 3$ (the peak).  Even as $F(t)$ decreases, friction alone cannot stop the block quickly.  The block continues moving well after the pulse ends.

    \item \textbf{Final displacement is an increasing function of $F_{\max}/(\mu_k mg)$.}  In the dimensionless form, the total displacement depends only on $\alpha$ and $\beta$ --- the ratios of friction forces to peak applied force.

    \item \textbf{If $\mu_s = \mu_k$} (no static/kinetic distinction), the onset time $t_1$ and the force-balance time $t_2$ become symmetric about $t = 3$.  The problem simplifies but the block still overshoots due to accumulated momentum.
\end{enumerate}

\end{document}
